\section*{Hoofdstuk \ref{ch:g11:subduction-zones} --- Subductiezones}

\begin{solution}{exo:g11:subduction-zones:1}
\omterm{def:g11:subduction-zones:subduction}{Subductie} is het wegzakken van een plaat \omterm{def:g11:ocean-ridges:oceanic-lithosphere}{oceanische lithosfeer} onder een andere
plaat, terug in de mantel; een \omterm{def:g11:subduction-zones:subduction}{subductiezone} is een convergente plaatgrens
waar dat gebeurt, gemarkeerd door een trog. \omterm{def:g11:ocean-ridges:oceanic-lithosphere}{Oceanische lithosfeer} zakt weg: zodra ze
oud is, is ze koud, dik en dichter dan de \omterm{def:g11:lithosphere-asthenosphere:lithosphere}{asthenosfeer} eronder. De
continentale korst is te licht om weg te zakken.
\end{solution}

\begin{solution}{exo:g11:subduction-zones:2}
(1) Bijvoorbeeld het gebied van Tonga en Fiji, het westen van de Grote Oceaan van Japan tot de
Marianen, Indonesië, en Zuid-Amerika. (2) Altijd aan de kant van de
bovenliggende plaat, nooit aan de oceaankant. (3) Ondiep: alle aardbevingen van
de Midden-Atlantische Rug op de kaart zijn ondieper dan \qty{70}{km}.
\end{solution}

\begin{solution}{exo:g11:subduction-zones:3}
$7031 \times 0.21 \approx 1480$ aardbevingen dieper dan \qty{70}{km};
$7031 \times 0.069 \approx 485$ dieper dan \qty{300}{km}.
\end{solution}

\begin{solution}{exo:g11:subduction-zones:4}
\omterm{def:g11:subduction-zones:eclogite}{Blauwschist} is een metamorf gesteente dat ontstaat uit het basalt van de \omterm{def:g11:ocean-ridges:oceanic-lithosphere}{oceanische korst}
dat tot hoge druk is begraven terwijl het relatief koud bleef, blauw gekleurd door
glaucofaan; \omterm{def:g11:subduction-zones:eclogite}{eclogiet} is het dichte metamorfe gesteente waarin basalt en \omterm{def:g11:ocean-ridges:gabbro}{gabbro}
op manteldiepte veranderen, opgebouwd uit rode granaat en groene pyroxeen. Beide komen uit
de \omterm{def:g11:ocean-ridges:oceanic-lithosphere}{oceanische korst} en behouden haar chemische samenstelling; wat verschilt, zijn de
mineralen, gevormd in vaste toestand onder de nieuwe druk en temperatuur,
en de dichtheid.
\end{solution}

\begin{solution}{exo:g11:subduction-zones:5}
Oceaanbodem: oude \omterm{def:g11:lithosphere-asthenosphere:lithosphere}{lithosfeer} die naar de trog beweegt. Trog: de plaat buigt
naar beneden en begint weg te zakken. Voorboog: sedimenten die van de wegzakkende plaat
zijn afgeschraapt, stapelen zich op; eronder schuiven de twee platen langs elkaar langs de plaatgrens,
waar zware aardbevingen plaatsvinden. Vulkanische boog: een rij vulkanen boven de
\omterm{def:g11:subduction-zones:slab}{slab}. Achterboog: achter sommige bogen wordt de bovenliggende plaat uitgerekt.
\end{solution}

\begin{solution}{exo:g11:subduction-zones:6}
$\tan\delta = (508 - 135)/(800 - 300) = 373/500 \approx 0.746$, dus
$\delta \approx \ang{37}$.
\end{solution}

\begin{solution}{exo:g11:subduction-zones:7}
(1) $\tan\ang{45} = 1$: \qty{300}{km} diep. (2) Op \qty{670}{km} van de
trog. (3) $670\sqrt{2} \approx \qty{948}{km}$ \omterm{def:g11:subduction-zones:slab}{slab}.
\end{solution}

\begin{solution}{exo:g11:subduction-zones:8}
(1) $(74.6 + 70.0 + 65.5 + 63.1 + 52.6)/5 \approx \qty{65.2}{mm}$ per jaar;
$(47.6 + 31.0 + 29.2 + 22.8 + 19.8 + 15.2 + 14.9 + 10.6)/8 \approx
\qty{23.9}{mm}$ per jaar; verhouding $65.2/23.9 \approx 2.7$. (2) Platen bewegen
het snelst wanneer er een \omterm{def:g11:subduction-zones:slab}{slab} aan hangt: de \omterm{def:g11:subduction-zones:slab-pull}{trek van de wegzakkende slab} is de
sterkste kracht die platen aandrijft, zoals Forsyth en Uyeda concludeerden.
\end{solution}

\begin{solution}{exo:g11:subduction-zones:9}
(1) \qty{83}{mm} is \qty{0.083}{m}: $50/0.083 \approx 600$ jaar en
$60/0.083 \approx 720$ jaar plaatbeweging. (2) Van 869 tot 2011 is 1142
jaar: van dezelfde orde. Als de plaatgrens tussen zware aardbevingen vastzit,
heeft de verschuiving van 2011 meerdere eeuwen opgeslagen beweging vrijgemaakt.
\end{solution}

\begin{solution}{exo:g11:subduction-zones:10}
(1) Japan: $6400 - 3200\,e^{-133/62.8} \approx \qty{6015}{m}$; Tonga:
$6400 - 3200\,e^{-99/62.8} \approx \qty{5740}{m}$; Andes:
$2500 + 350\sqrt{49} = \qty{4950}{m}$. (2) De bodem die Japan bereikt: de
oudste, dus de koudste en dichtste \omterm{def:g11:lithosphere-asthenosphere:lithosphere}{lithosfeer}, en de diepste bodem.
\end{solution}

\begin{solution}{exo:g11:subduction-zones:11}
(1) $3380 \times 1.06 \approx \qty{3583}{kg/m^3}$. (2)
$(3583 - 2900)/2900 \approx 0.24$: een stijging van ongeveer \qty{24}{\percent}.
(3) Ervoor is de korst lichter dan de mantel met
$3380 - 2900 = \qty{480}{kg/m^3}$; erna is ze dichter met ongeveer
\qty{203}{kg/m^3}. Ze houdt de \omterm{def:g11:subduction-zones:slab}{slab} niet meer omhoog: ze trekt hem naar beneden.
\end{solution}

\begin{solution}{exo:g11:subduction-zones:12}
(1) $t = 6 \times 10^6 \times 3.16 \times 10^7 \approx 1.9 \times 10^{14}$~s;
$\sqrt{10^{-6} \times 1.9 \times 10^{14}} \approx \qty{1.4e4}{m}$, ongeveer
\qty{14}{km}. (2) Veel minder dan \qty{100}{km}: in 6 miljoen jaar warmt alleen de
buitenste huid van de \omterm{def:g11:subduction-zones:slab}{slab} op; zijn binnenste blijft koud. (3) Koud gesteente is
star en kan breken, terwijl de hete mantel eromheen vloeit; de diepe
aardbevingen vinden plaats in de koude \omterm{def:g11:subduction-zones:slab}{slab}.
\end{solution}

\begin{solution}{exo:g11:subduction-zones:13}
(1) $\sqrt{1017^2 + 670^2} \approx \qty{1218}{km}$. (2) \qty{64}{mm} per jaar
is \qty{64}{km} per miljoen jaar: $1218/64 \approx 19$ miljoen jaar.
(3) Ongeveer drie keer zo lang als onder Tonga: de \omterm{def:g11:subduction-zones:slab}{slab} onder de Andes
zakt langzamer en met een flauwere helling.
\end{solution}

\begin{solution}{exo:g11:subduction-zones:14}
(1) Een plak van \qty{1}{m} breed langs de trog heeft een volume van
$10^5 \times 10^6 \times 1 = \qty{1e11}{m^3}$; zijn massaoverschot is
$36 \times 10^{11} = \qty{3.6e12}{kg}$. (2) Zijn overgewicht is
$3.6 \times 10^{12} \times 9.8 \approx \qty{3.5e13}{N}$ per meter
trog. (3) De \omterm{def:g11:subduction-zones:slab}{slab} zit vast aan de rest van de plaat, die star is:
de kracht wordt overgebracht op de hele plaat, die naar de trog wordt
gesleept, des te meer omdat troggen duizenden kilometers lang zijn.
\end{solution}

\begin{solution}{exo:g11:subduction-zones:15}
(1) $(23.6 - 16.5) \times 111 \approx \qty{790}{km}$. (2) De snelheid neemt toe met
$237 - 104 = \qty{133}{mm}$ per jaar over \qty{790}{km}, ongeveer \qty{17}{mm} per
jaar per honderd kilometer. (3) De convergentie is de snelheid van de
Pacifische plaat min die van de boog. Omdat de snelheid van de Pacifische plaat op
beide punten ongeveer gelijk is, moet de boog zelf naar de
Pacifische plaat toe bewegen, steeds sneller naar het noorden: de achterboog erachter wordt
uitgerekt, meer in het noorden dan in het zuiden.
\end{solution}

\begin{solution}{pb:g11:subduction-zones:1}
\textbf{1.} Ongeveer \qty{671}{km}.

\textbf{2.} Aan de oceaankant is de plaat nog niet weggezakt: de \omterm{def:g11:lithosphere-asthenosphere:lithosphere}{lithosfeer}
ligt daar dicht bij het oppervlak. De diepe aardbevingen vinden plaats in de \omterm{def:g11:subduction-zones:slab}{slab}, die
alleen onder de bovenliggende plaat ligt.

\textbf{3.} $z(300) = -24 + 273 = \qty{249}{km}$ en
$z(700) = -24 + 637 = \qty{613}{km}$; $\tan\delta = 364/400 = 0.91$, dus
$\delta \approx \ang{42}$.

\textbf{4.} $x = (100 + 24)/0.91 \approx \qty{136}{km}$ en
$x = (400 + 24)/0.91 \approx \qty{466}{km}$.

\textbf{5.} Een rechte lijn met deze helling die bij de trog begint, zou
door $z = 0$ gaan bij $x = 0$; de lijn ligt erboven. De \omterm{def:g11:subduction-zones:slab}{slab} is niet recht: dicht bij de
trog helt hij flauwer, en dieper wordt hij steiler.

\textbf{6.} \qty{16.8}{cm} per jaar, \qty{168}{km} per miljoen jaar.

\textbf{7.} $168/69 \approx 2.4$ keer sneller.

\textbf{8.} De convergentie over de trog is de snelheid van de Pacifische
plaat min de snelheid van de boog aan de andere kant. Als de boog naar
de Pacifische plaat toe beweegt, zoals hij doet terwijl de achterboog erachter wordt uitgerekt, tellen de
twee snelheden op, en is de convergentie groter dan de snelheid van de
Pacifische plaat alleen.

\textbf{9.} $6400 - 3200\,e^{-99/62.8} \approx \qty{5740}{m}$.

\textbf{10.} Na ongeveer 60 miljoen jaar is \omterm{def:g11:ocean-ridges:oceanic-lithosphere}{oceanische lithosfeer} als geheel dichter
dan de \omterm{def:g11:lithosphere-asthenosphere:lithosphere}{asthenosfeer} eronder; bij 99 miljoen jaar is ze nog dichter.
Zodra haar rand naar beneden gaat, trekt haar gewicht de rest erin.

\textbf{11.} $(670 + 24)/0.91 \approx \qty{763}{km}$.

\textbf{12.} $\sqrt{763^2 + 670^2} \approx \qty{1015}{km}$.

\textbf{13.} $1015/168 \approx 6.0$ miljoen jaar.

\textbf{14.} $1015/69 \approx 14.7$ miljoen jaar.

\textbf{15.} $3380 \times 1.06 \approx \qty{3583}{kg/m^3}$: de korst, eerst
lichter dan de mantel, wordt er dichter dan, en de hele \omterm{def:g11:subduction-zones:slab}{slab} trekt
harder.

\textbf{16.} $t \approx 6.0 \times 10^6 \times 3.16 \times 10^7 \approx
1.9 \times 10^{14}$~s, en $\sqrt{\kappa t} \approx \qty{14}{km}$.

\textbf{17.} \qty{14}{km} is weinig naast \qty{100}{km}: alleen het buitenste deel
van de \omterm{def:g11:subduction-zones:slab}{slab} is opgewarmd; zijn binnenste is nog koud wanneer hij
\qty{670}{km} bereikt.

\textbf{18.} De koude \omterm{def:g11:subduction-zones:slab}{slab} blijft star genoeg om te breken, tot aan zijn punt,
terwijl de hete mantel eromheen vloeit: de aardbevingen onder Tonga vinden plaats
in de \omterm{def:g11:subduction-zones:slab}{slab}, tot ongeveer \qty{671}{km}.

\textbf{19.} Seismische golven tonen een scherpe verandering in de mantel op ongeveer
\qty{670}{km}: hun snelheden springen daar omhoog.

\textbf{20.} Ongeveer \textbf{6 miljoen jaar} van de trog tot
\qty{670}{km}.
\end{solution}
